\documentclass[12pt]{article}

% -------- PACKAGES ---------
\usepackage{cosc343style}
\usepackage{subfig}
\usepackage{lipsum}
\usepackage{bbm}
\usepackage{graphicx}
\usepackage{listings}
\usepackage{xcolor}
\usepackage{url}
\usepackage{eurosym}
\usepackage{amsmath}
\usepackage{amssymb}
\usepackage{soul}
\usepackage{xcolor}
\usepackage{subcaption}
\usepackage{float}
\sethlcolor{yellow!80}
\usepackage{algorithm}
\usepackage{algpseudocode}
\usepackage{amsthm}
\usepackage{hyperref}

% -------- COMMANDS ---------
\newcommand*{\N}{\mathbb N}
\newcommand*{\Z}{\mathbb Z}
\newcommand*{\Q}{\mathbb Q}
\newcommand*{\R}{\mathbb R}
\newcommand*{\C}{\mathbb C}
\newcommand{\B}{\mathcal B}
\newcommand{\A}{\mathcal A}
\newcommand{\ip}[2]{\left\langle #1 \middle| #2 \right\rangle}
\newcommand{\norm}[1]{\left\lVert #1 \right\rVert}
\newcommand{\sech}{\mathrm{sech}}
\newcommand{\csch}{\mathrm{csch}}
\newcommand*{\e}{\mathrm{e}}
\newcommand*{\im}{\mathrm{i}}
\newcommand*{\cint}{\oint\limits}
\newcommand*{\rmd}{\mathrm{d}}
\newcommand*{\D}{\mathcal{D}}
\newcommand{\Pm}{\textbf{P}}
\newcommand{\ind}{\mathbbm{1}} 
\newcommand{\Ell}{\mathcal{L}}
\newcommand{\F}{\mathcal F}
\newcommand*{\E}{\mathcal{E}}

\makeatletter
\@ifundefined{href}{%
  \protected\def\href#1#2{\url{#1}}%
}{}
\makeatother

\definecolor{burgundy}{RGB}{200,0,32}

\makeatletter
\renewenvironment{thebibliography}[1]{%
  \par\bigskip
  {\large \bfseries References\par}\medskip
  \begin{list}{[\arabic{enumiv}]}{%
    \usecounter{enumiv}%
    \setlength{\labelsep}{0.5em}%
    \setlength{\leftmargin}{2em}%
    \setlength{\itemsep}{0.25\baselineskip}%
  }%
  \sloppy\clubpenalty4000\widowpenalty4000\sfcode`\.=1000%
}{%
  \end{list}%
}
\makeatother

\lstset{
  language=R,
  basicstyle=\ttfamily\scriptsize,
  commentstyle=\color{green!40!black},
  stringstyle=\color{burgundy},
  numbers=left,
  numberstyle=\tiny\color{gray},
  stepnumber=1,
  numbersep=8pt,
  breaklines=true,
  showstringspaces=false,
  backgroundcolor=\color{gray!10}, 
}

\lstdefinestyle{Routput}{
  basicstyle=\ttfamily\scriptsize,
  numbers=none,
  breaklines=true,
  showstringspaces=false, 
  backgroundcolor=\color{white}, 
  commentstyle=\color{black},
  stringstyle=\color{black},
}

\lstdefinestyle{Routput2}{
  basicstyle=\ttfamily\fontsize{6}{6}\selectfont, 
  numbers=none,
  breaklines=true,
  showstringspaces=false, 
  backgroundcolor=\color{white}, 
  commentstyle=\color{black},
  stringstyle=\color{black},
}

\newtheorem{theorem}{Theorem}[section]
\newtheorem{lemma}[theorem]{Lemma}
\newtheorem{proposition}[theorem]{Proposition}
\newtheorem{corollary}[theorem]{Corollary}

\theoremstyle{definition}
\newtheorem{definition}[theorem]{Definition}

\theoremstyle{remark}
\newtheorem{remark}[theorem]{Remark}
\newtheorem*{remark*}{Remark}

\theoremstyle{plain}
\newtheorem{example}[theorem]{Example}


% -------- HEADER ---------
\papercode{Rat-Mice}
\title{Ergodicity}
\author{Stella \textsc{Srzich}}
\studentid{3371930}
\reportdate{\today}


% -------- DOCUMENT ---------
\allowdisplaybreaks
\begin{document}

\maketitle

% \begin{lstlisting}[language=Python]
%     example
% \end{lstlisting} 
% \begin{lstlisting}[style=Routput]
%     example
% \end{lstlisting} 

\textit{This document is meant to serve as a 
discussion of the Rat-Mice algorithm and sketch of the proof of its ergodicity.}


\section{What kind of algorithm do we want?}

In general, we want an algorithm that 
correctly samples from a target density $\pi(\cdot)$ 
(i.e., produces a Markov chain that is ergodic with respect to $\pi(\cdot)$). 
Moreover, we would like the chain to converge to the target density within a reasonable number of iterations. 
Moreover, out of all such algorithms, we would like to minimize the computational cost per effective sample.

The standard algorithm that indeed correctly samples 
from a target density $\pi(\cdot)$ is the Metropolis--Hastings (MH) algorithm.

\begin{algorithm}[H]
  \caption{Metropolis--Hastings (MH)}
  \label{alg:mh}

  \textbf{function:} $[x^1,\ldots,x^N] =
  \mathrm{MH}\bigl(\pi(\cdot), q(\cdot\mid\cdot), x^0, N\bigr)$

  \textbf{input:} target density $\pi(\cdot)$, proposal density
  $q(\cdot\mid\cdot)$, initial state $x^0$, number of steps $N$

  \textbf{output:} ordered list of states $[x^1,\ldots,x^N]$

  \begin{algorithmic}[1]
  \For{$j = 0$ to $N-1$}
      \State Given $x^j$, generate a proposal $y$ distributed as $q(y \mid x^j)$.
      \State Accept proposal $y$ as the next state, i.e. set $x^{j+1} = y$, with probability
      \[
          \alpha(y \mid x^j)
          =
          \min\left\{
              1,
              \frac{\pi(y)q(x^j\mid y)}
                   {\pi(x^j)q(y\mid x^j)}
          \right\}.
      \]
      \State Otherwise reject $y$ and set $x^{j+1} = x^j$.
  \EndFor
  \end{algorithmic}
\end{algorithm}

However, when the target density $\pi(\cdot)$ is multimodal 
(e.g., as is the case when we use high-level representations for states), 
with basic choices of proposal kernel $q(\cdot\mid\cdot)$,
the MH algorithm can struggle to explore all modes within a reasonable number of iterations.
In this case, the MH algorithm will not converge to the target density within a reasonable number of iterations. 

One way to address this is to, at each iteration, run a subchain of length $n$ that targets 
a coarse (i.e., cheap to evaluate) density $\pi_{C}(\cdot)$, and then use the resulting state as a proposal for this iteration.
The idea is that 
if the coarse density $\pi_{C}(\cdot)$ approximates the target density $\pi(\cdot)$ well enough,
after a sufficient number of subchain steps, the proposal will be in a different mode. 
This effectively produces a new proposal kernel (from the basic proposal kernel $q(\cdot\mid\cdot)$) for use in 
in the MH algorithm that ``wormholes'' between modes without us having to be clever about explicitly designing such a proposal. 
Or, as we like to think of it, a proposal kernel of little ``rats and mice'' that tunnel between the modes.
This idea is implemented in the Randomised-Length-Subchain Surrogate Transition (RST) algorithm. 

\begin{algorithm}[H]
  \caption{Randomised-Length-Subchain Surrogate Transition (RST)}
  \label{alg:rst}
  \textbf{function:} $[x^1,\ldots,x^N] =
  \mathrm{RST}\bigl(
      \pi_F(\cdot),
      \pi_C(\cdot),
      q(\cdot\mid\cdot),
      p(\cdot),
      x^0,
      N
  \bigr)$
  
  \textbf{input:} target (fine) density $\pi_F(\cdot)$, surrogate (coarse) density $\pi_C(\cdot)$,
  proposal kernel $q(\cdot\mid\cdot)$, probability mass function $p(\cdot)$ over subchain length,
  initial state $x^0$, number of steps $N$
  
  \textbf{output:} ordered list of states $[x^1,\ldots,x^N]$
  (or just the final state $x^N$)
  
  \begin{algorithmic}[1]
  \For{$j = 0$ to $N-1$}
      \State Draw the subchain length $n \sim p(\cdot)$.

      \State Starting at $x^j$, generate a subchain of length $n$ using MH Alg.~1 to target $\pi_C(\cdot)$:
      \[
          y
          =
          \mathrm{MH}\bigl(
              \pi_C(\cdot),
              q(\cdot\mid\cdot),
              x^j,
              n
          \bigr).
      \]
      
      \State Accept the proposal $y$ as the next sample, i.e. set
      $x^{j+1} = y$, with probability
      \[
          \alpha(y \mid x^j)
          =
          \min\left\{
              1,
              \frac{
                  \pi_F(y)\pi_C(x^j)
              }{
                  \pi_F(x^j)\pi_C(y)
              }
          \right\}.
      \]
      
      \State Otherwise reject $y$ and set $x^{j+1} = x^j$.
  \EndFor
  \end{algorithmic}
\end{algorithm}

To make the RST algorithm more efficient, 
one would like to estimate the mean and covariance of the discrepancy between the exact and coarse forward maps 
and subsequently modify the coarse likelihood so that the resulting coarse density more closely approximates the target density as the chain progresses
(see Fox, Cui and Neumayer 2020).

If we allow state-dependent coarse densities $\pi_{C,x}(\cdot)$,
as opposed to a single state-independent coarse density $\pi_{C}(\cdot)$ as in the RST algorithm,
the coarse forward map can be locally corrected so that it agrees with the exact forward map at the current state $x$.
For this locally corrected map, the discrepancy is zero at $x$ and has mean zero under the stationary transition law. 
For some reason (I think because we no longer need to correct the global coarse model error, we only need to correct the 
change in the coarse model error over one move), the modified
state-dependent coarse densities more closely approximate the target density than the modified state-independent coarse density.
Long story short, this results in a more efficient algorithm.

We name this algorithm the Rat-Mice (RM) algorithm.
For simplicity, we begin by only considering a fixed subchain length $n$.

\begin{algorithm}[H]
  \caption{Rat-Mice (RM)}
  \label{alg:rm}
  \textbf{function:} $[x^1,\ldots,x^N] =
  \mathrm{RM}\bigl(\pi(\cdot), \{\pi_{C,x}(\cdot)\}_{x \in \mathcal{X}}, q(\cdot\mid\cdot), n, x^0, N\bigr)$
  
  \textbf{input:} target density $\pi(\cdot)$, state-dependent surrogate (coarse) densities $\{\pi_{C,x}(\cdot)\}_{x \in \mathcal{X}}$, 
  proposal kernel $q(\cdot\mid\cdot)$, subchain length $n$, initial state $x^0$, number of steps $N$
  
  \textbf{output:} ordered list of states $[x^1,\ldots,x^N]$
  
  \begin{algorithmic}[1]
  \For{$j = 0$ to $N-1$}
      \State Starting at $x^j$, generate subchain of length $n$ using MH Alg. 1 to target $\pi_{C, x_j}(\cdot)$:
      \[
          y = \mathrm{MH}\bigl(\pi_{C, x_j}(\cdot), q(\cdot\mid\cdot), x^j, n\bigr).
      \]
      \State Accept the state $y$ as our proposal, i.e. set $x^{j+1} = y$, with probability
      \[
          \alpha(y \mid x^j)
          =
          \text{(some value that ensures the chain is in detailed balance with $\pi(\cdot)$)}.
      \]
      \State Otherwise reject $y$ and set $x^{j+1} = x^j$.
  \EndFor
  \end{algorithmic}
\end{algorithm}

We now recap the point of the RM algorithm. 
If you use basic proposal kernels (i.e., ones you are not super clever about), 
the MH algorithm will not converge to a multimodal target density within a reasonable number of iterations.
One way to address this is to use the basic proposal kernel to generate a more 
sophisticated proposal kernel that acts like little ``rats and mice'' that tunnel between the modes. 

To generate this more sophisticated proposal kernel, you want to target a density that is cheap to evaluate 
(so that doing a bunch of steps in the subchain is cheap), and you want the density to be a good approximation 
of the target density (so that proposals kind of look like samples from the target density and are accepted with high probability). 
The RST algorithm does this, but only allows for a single state-independent coarse density.

Given a state-independent coarse density (which will be defined by a coarse forward map), we can easily 
define state-dependent coarse densities that are better approximations of the target density 
(and are as cheap to evaluate as they use the same coarse forward map).
Thus, 
we want to somehow modify the RST algorithm to allow for state-dependent coarse densities.


\section{An ideal world}

\begin{definition}[Peskun ordering]
Let $P_1$ and $P_2$ be Markov transition kernels on
$(\mathcal X,\mathcal B(\mathcal X))$ with common invariant distribution
$\pi$. We say that $P_1$ \emph{dominates} $P_2$ in the Peskun ordering,
and write
\[
P_1 \succeq_{\mathrm P} P_2,
\]
if, for every $x\in\mathcal X$ and every measurable set
$A\in\mathcal B(\mathcal X)$ such that $x\notin A$,
\[
P_1(x,A)\geq P_2(x,A).
\]
\end{definition}

\begin{theorem}[Peskun ordering and asymptotic variance]
Let $P_1$ and $P_2$ 
be Markov transition kernels that
satisfy detailed balance with respect to $\pi(\cdot)$ with
\[
P_1 \succeq_{\mathrm P} P_2.
\]
Then, for every $f\in L^2(\pi)$ for which the relevant asymptotic
variances are well defined,
\[
\sigma_{P_1}^2(f)
\leq
\sigma_{P_2}^2(f),
\]
where
\[
\sigma_P^2(f)
:=
\lim_{N\to\infty}
N\,
\operatorname{Var}_{\pi}\left(
\frac{1}{N}\sum_{j=1}^N f(X_j)
\right)
\]
for a stationary Markov chain with transition kernel $P$.
\end{theorem}

\begin{theorem}[Pointwise maximality of the Metropolis--Hastings acceptance probability]
Let $\pi(\cdot)$ be a target density on $\mathcal X$, and let $q(\cdot\mid\cdot)$ be a
proposal density. Suppose that an acceptance probability
$\alpha:\mathcal X\times\mathcal X\to[0,1]$ satisfies the detailed balance
condition
\[
\pi(x)q(y\mid x)\alpha(x,y)
=
\pi(y)q(x\mid y)\alpha(y,x)
\]
for all $x,y\in\mathcal X$.
Then, for all $x,y\in\mathcal X$ with $\pi(x)>0$ and $q(y\mid x)>0$,
\[
\alpha(x,y)
\leq
\alpha_{\mathrm{MH}}(x,y)
:=
\min\left\{
1,
\frac{\pi(y)q(x\mid y)}
      {\pi(x)q(y\mid x)}
\right\}.
\]
\end{theorem}

\begin{proof}
Fix $x,y\in\mathcal X$ with $\pi(x)>0$ and $q(y\mid x)>0$.
Moreover, detailed balance gives
\[
\alpha(x,y)
=
\frac{\pi(y)q(x\mid y)}
      {\pi(x)q(y\mid x)}
\alpha(y,x).
\]
Since $\alpha(y,x)\leq 1$, it follows that
\[
\alpha(x,y)
\leq
\frac{\pi(y)q(x\mid y)}
      {\pi(x)q(y\mid x)}.
\]
Combining the two bounds yields
\[
\alpha(x,y)
\leq
\min\left\{
1,
\frac{\pi(y)q(x\mid y)}
      {\pi(x)q(y\mid x)}
\right\}
=
\alpha_{\mathrm{MH}}(x,y).
\]
\end{proof}

\begin{theorem}[Optimality of the Metropolis--Hastings acceptance probability]
Fix a target density $\pi(\cdot)$ and a proposal density $q(\cdot\mid\cdot)$. Let
$P_{\mathrm{MH}}$ be the Metropolis--Hastings transition kernel with
acceptance probability $\alpha_{\mathrm{MH}}$, and let $P_{\alpha}$ be
any other accept--reject transition kernel using the
same proposal $q(\cdot\mid\cdot)$ and an acceptance probability $\alpha(\cdot\mid\cdot)$ satisfying
detailed balance with respect to $\pi(\cdot)$.
Then
\[
P_{\mathrm{MH}}\succeq_{\mathrm P}P_{\alpha}.
\]
Consequently, for every $f\in L^2(\pi)$ for which the relevant
asymptotic variances are well defined,
\[
\sigma_{P_{\mathrm{MH}}}^2(f)
\leq
\sigma_{P_{\alpha}}^2(f).
\]
Thus, among all accept--reject kernels using the fixed
proposal $q$ and an acceptance probability $\alpha(\cdot\mid\cdot)$ 
satisfying detailed balance with respect to $\pi(\cdot)$, the Metropolis--Hastings acceptance
probability is optimal in the sense of asymptotic variance.
\end{theorem}

\begin{proof}
For every measurable set $A\in\mathcal B(\mathcal X)$ such that
$x\notin A$,
\[
P_{\alpha}(x,A)
=
\int_A q(y\mid x)\alpha(x,y)\,dy
\]
and
\[
P_{\mathrm{MH}}(x,A)
=
\int_A q(y\mid x)\alpha_{\mathrm{MH}}(x,y)\,dy.
\]
By the pointwise maximality of the Metropolis--Hastings acceptance
probability,
\[
q(y\mid x)\alpha_{\mathrm{MH}}(x,y)
\geq
q(y\mid x)\alpha(x,y).
\]
Therefore,
\[
P_{\mathrm{MH}}(x,A)
\geq
P_{\alpha}(x,A),
\]
so
\[
P_{\mathrm{MH}}\succeq_{\mathrm P}P_{\alpha}.
\]
The asymptotic-variance inequality now follows from the Peskun-ordering
theorem.
\end{proof}

From the above theorem, we can see that for the fixed effective proposal kernel 
$q'(y\mid x) = q_x^n(y\mid x)$, 
if we want to choose an acceptance probability $\alpha(\cdot\mid\cdot)$ so that 
the resulting transition kernel satisfies detailed balance with respect to $\pi(\cdot)$,
then the MH acceptance probability is the optimal acceptance probability in the sense of asymptotic variance.
This ideal Rat-Mice (IRM) algorithm is given below.

\begin{algorithm}[H]
  \caption{Ideal Rat-Mice (IRM)}
  \label{alg:irm}
  \textbf{function:} $[x^1,\ldots,x^N] =
  \mathrm{IRM}\bigl(\pi(\cdot), \{\pi_{C,x}(\cdot)\}_{x \in \mathcal{X}}, q(\cdot\mid\cdot), n, x^0, N\bigr)$
  
  \textbf{input:} target density $\pi(\cdot)$, state-dependent surrogate (coarse) densities $\{\pi_{C,x}(\cdot)\}_{x \in \mathcal{X}}$, 
  proposal kernel $q(\cdot\mid\cdot)$, subchain length $n$, initial state $x^0$, number of steps $N$
  
  \textbf{output:} ordered list of states $[x^1,\ldots,x^N]$
  
  \begin{algorithmic}[1]
  \For{$j = 0$ to $N-1$}
      \State Starting at $x^j$, generate subchain of length $n$ using MH Alg. 1 to target $\pi_{C, x_j}(\cdot)$:
      \[
          y = \mathrm{MH}\bigl(\pi_{C, x_j}(\cdot), q(\cdot\mid\cdot), x^j, n\bigr).
      \]
      \State Accept the state $y$ as our proposal, i.e. set $x^{j+1} = y$, with probability
      \[
          \alpha(y \mid x^j)
          =
          \min\left\{1, \frac{\pi(y)q_y^n(x^j|y)}{\pi(x^j)q_{x^j}^n(y|x^j)}\right\},
      \]
      where 
      \[
        q_x(b|a) = q(b|a) \min\left\{1, \frac{\pi_{C,x}(b)q(a|b)}{\pi_{C,x}(a)q(b|a)}\right\} + \delta_a(b)\left(1 - \int_\mathcal X q(c|a) \min\left\{1, \frac{\pi_{C,x}(c)q(a|c)}{\pi_{C,x}(a)q(c|a)}\right\} dc\right).
      \]
      \State Otherwise reject $y$ and set $x^{j+1} = x^j$.
  \EndFor
  \end{algorithmic}
\end{algorithm}

\begin{remark}
  Let us recall what it means for an acceptance probability to be optimal in the sense of asymptotic variance.
  It means that the asymptotic variance appearing in the central limit theorem for the ergodic average is minimized.
  That is, for a large enough fixed number of iterations, the ergodic average will have smaller Monte Carlo variance than any other acceptance probability. 
  This does not necessarily imply faster convergence of the chain to its invariant distribution. 

  That is, if the fixed number of iterations is not large enough, 
  the ergodic average of the chain need not be distributed according to the central limit theorem, 
  and so need not probably be closer to the true mean than the ergodic average of the chain using any other acceptance probability.

  All of this is to say that if the chain takes ages to converge, 
  the MH acceptance probability may not be the optimal choice. 
  With that said, we hope that the Rat-Mice effective proposal 
  will make the chain converge fast enough so that the ergodic average 
  is distributed according to the central limit theorem in a reasonable amount of time, 
  and the MH acceptance probability is the optimal choice.
  Moreover, designing algorithms that converge faster is a whole nother can of worms.
\end{remark}

\section{What choice of coarse density is best for IRM?}

There are two competing intuitions about what choice of coarse density is best.
One is that you want the coarse density to be close to target so that the proposal from the subchain 
looks like its from the target density and hopefully is in one of its other modes. 
The other is that you want it to be flatter than the target density so that you can actually go through the 
regions of low probability to the modes. 

My initial thoughts where that if you have two (for now, state-independent) 
coarse densities $\pi_{C,1}(\cdot)$ and $\pi_{C,2}(\cdot)$, and one dominates the other 
in the DA acceptance probability sense, i.e., 
\begin{align*}
  \min\left\{1, \frac{\pi(y)\pi_{C,1}(x)}{\pi(x)\pi_{C,1}(y)}\right\}
  \geq
  \min\left\{1, \frac{\pi(y)\pi_{C,2}(x)}{\pi(x)\pi_{C,2}(y)}\right\}
\end{align*}
for all $x, y\in\mathcal X$ with $\pi(x)>0$ and $q(y\mid x)>0$,
then the algorithm using $\pi_{C,1}(\cdot)$ would Peskun-dominate the algorithm using $\pi_{C,2}(\cdot)$.
Thus, the algorithm using $\pi_{C,1}(\cdot)$ would be the optimal in the sense of asymptotic variance.

However, it turns out that this is not the case. 
Indeed, if $\pi_{C,1}(\cdot)$ dominates $\pi_{C,2}(\cdot)$ in the DA acceptance probability sense, 
since the effective proposals will be in detailed balance with respect to $\pi_{C,1}(\cdot)$ and $\pi_{C,2}(\cdot)$ respectively, 
the outer acceptance probability is precisely the DA acceptance probability. 
Thus, indeed, if $\pi_{C,1}(\cdot)$ dominates $\pi_{C,2}(\cdot)$ in the DA acceptance probability sense, 
the algorithm using $\pi_{C,1}(\cdot)$ will accept more proposals than the algorithm using $\pi_{C,2}(\cdot)$.

However, Peskun ordering also depends on the effective proposal. 
And, the effective proposal is different for the two algorithms. 
So, the domination of the outer acceptance probability does not imply peskun ordering.
In fact, we provide a counter example below. 

\begin{example}[DA-domination does not imply Peskun-domination]
Let
\[
\mathcal X=\{1,2,3\},
\qquad
\pi=\left(\frac13,\frac13,\frac13\right),
\]
and use the same proposal for both inner chains:
\[
q(y\mid x)=\frac12,
\qquad y\neq x.
\]
Let
\[
\pi_1=(0.22,0.43,0.35),
\qquad
\pi_2=(0.01,0.63,0.36),
\]
and let \(Q_i\) denote the Metropolis--Hastings kernel with target
\(\pi_i\) and proposal \(q\). Take the inner-chain length to be \(n=2\).

Since \(Q_i^2\) is reversible with respect to \(\pi_i\), and \(\pi\) is
uniform, the outer acceptance probability is
\[
\alpha_i(x,y)
=
\min\left\{1,\frac{\pi_i(x)}{\pi_i(y)}\right\}.
\]
Both coarse densities satisfy
\[
\pi_i(2)>\pi_i(3)>\pi_i(1).
\]
For the directions in which the outer acceptance probability is less
than one,
\[
\alpha_1(1,2)
=
\frac{0.22}{0.43}
>
\frac{0.01}{0.63}
=
\alpha_2(1,2),
\]
\[
\alpha_1(1,3)
=
\frac{0.22}{0.35}
>
\frac{0.01}{0.36}
=
\alpha_2(1,3),
\]
and
\[
\alpha_1(3,2)
=
\frac{0.35}{0.43}
>
\frac{0.36}{0.63}
=
\alpha_2(3,2).
\]
In the reverse directions, both acceptance probabilities equal one.
Hence
\[
\alpha_1(x,y)\geq\alpha_2(x,y)
\qquad
\text{for every }x\neq y.
\]

However, direct calculation gives
\[
Q_1^2(3\mid2)
=
\frac{630}{1849}
\approx 0.34072,
\]
whereas
\[
Q_2^2(3\mid2)
=
\frac{152}{441}
\approx 0.34467.
\]
Moreover,
\[
\alpha_1(2,3)=\alpha_2(2,3)=1.
\]
Therefore,
\[
Q_1^2(3\mid2)\alpha_1(2,3)
<
Q_2^2(3\mid2)\alpha_2(2,3).
\]
Thus the first algorithm does not Peskun-dominate the second, despite
having a pointwise larger outer acceptance probability.
\end{example}

Moreover, if we allow the coarse densities to be state-dependent, 
dominance in the DA acceptance probability sense does not even imply dominance in the outer acceptance probability sense.
See the counter example below. 

\begin{example}[Statewise DA-domination does not imply outer-acceptance domination]
Consider
\[
\mathcal X=\{0,1\},
\qquad
\pi(0)=\pi(1)=\frac12,
\]
and take the symmetric deterministic proposal
\[
q(1\mid0)=q(0\mid1)=1.
\]
Let the inner-chain length be $n=1$.

For algorithm $1$, choose
\[
\pi_{C,1,0}
=
\left(\frac34,\frac14\right),
\qquad
\pi_{C,1,1}
=
\left(\frac1{10},\frac9{10}\right),
\]
where the entries are ordered according to the states $(0,1)$.

For algorithm $2$, choose
\[
\pi_{C,2,0}
=
\left(\frac15,\frac45\right),
\qquad
\pi_{C,2,1}
=
\left(\frac3{10},\frac7{10}\right).
\]

Since $\pi$ is uniform,
\[
\beta_i(x,y)
:=
\min\left\{
1,
\frac{\pi(y)}{\pi(x)}
\frac{\pi_{C,i,y}(x)}
      {\pi_{C,i,x}(y)}
\right\}
=
\min\left\{
1,
\frac{\pi_{C,i,y}(x)}
      {\pi_{C,i,x}(y)}
\right\}.
\]
For the move $0\to1$,
\[
\beta_1(0,1)
=
\frac{1/10}{1/4}
=
\frac25,
\]
whereas
\[
\beta_2(0,1)
=
\frac{3/10}{4/5}
=
\frac38.
\]
Thus,
\[
\beta_1(0,1)>\beta_2(0,1).
\]

For the reverse move,
\[
\beta_1(1,0)
=
\min\left\{1,\frac{1/4}{1/10}\right\}
=
1,
\]
and
\[
\beta_2(1,0)
=
\min\left\{1,\frac{4/5}{3/10}\right\}
=
1.
\]
Hence
\[
\beta_1(x,y)\geq\beta_2(x,y)
\qquad
\text{for every }x\neq y.
\]

Let $Q_{i,x}$ denote the one-step inner Metropolis--Hastings kernel
targeting $\pi_{C,i,x}$. For algorithm $1$,
\[
Q_{1,0}(1\mid0)
=
\min\left\{1,\frac{1/4}{3/4}\right\}
=
\frac13,
\]
and
\[
Q_{1,1}(0\mid1)
=
\min\left\{1,\frac{1/10}{9/10}\right\}
=
\frac19.
\]
Therefore,
\[
\alpha_1^{\mathrm{out}}(0,1)
=
\min\left\{
1,
\frac{\pi(1)Q_{1,1}(0\mid1)}
      {\pi(0)Q_{1,0}(1\mid0)}
\right\}
=
\frac{1/9}{1/3}
=
\frac13.
\]

For algorithm $2$,
\[
Q_{2,0}(1\mid0)
=
\min\left\{1,\frac{4/5}{1/5}\right\}
=
1,
\]
and
\[
Q_{2,1}(0\mid1)
=
\min\left\{1,\frac{3/10}{7/10}\right\}
=
\frac37.
\]
Thus,
\[
\alpha_2^{\mathrm{out}}(0,1)
=
\min\left\{
1,
\frac{\pi(1)Q_{2,1}(0\mid1)}
      {\pi(0)Q_{2,0}(1\mid0)}
\right\}
=
\frac37.
\]

Consequently,
\[
\alpha_1^{\mathrm{out}}(0,1)
=
\frac13
<
\frac37
=
\alpha_2^{\mathrm{out}}(0,1),
\]
despite algorithm $1$ dominating algorithm $2$ in the stated
DA sense.
\end{example}

Overall, we have seen that for state-independent coarse densities, 
domination in the DA acceptance probability sense 
implies domination in the outer acceptance probability sense, but not 
in the Peskun ordering sense.
For state-dependent coarse densities, we have seen that domination in the DA acceptance probability sense 
does not even imply domination in the outer acceptance probability sense.

This means that we do not necessarily want the coarse density to be close to the target density 
in the DA acceptance probability sense. 
Instead, we want the coarse density to be close to the target density in the 
sense of dominating in terms of the following quantity
\begin{align*}
  q_{x}^{n}(y\mid x)\alpha(x,y)
  =
  \min\left\{
  q_{x}^{n}(y\mid x),
  \frac{\pi(y)}{\pi(x)}q_{y}^{n}(x\mid y)
  \right\},
\end{align*}
which is equivalent to domination in the Peskun ordering sense for the IRM algorithm.

With that said, I still think that driving towards a coarse density 
that is cheap to evaluate and close to the target is a reasonable thing to do.
And, in any case, the choice of coarse density should not 
affect the general structure of the RM algorithm I develop.


\section{How does IRM compare to DA?}

\begin{lemma}[Acceptance probability for DA with state-dependent coarse densities]
Fix a target density $\pi(\cdot)$, state-dependent coarse densities $\{\pi_{C,x}(\cdot)\}_{x \in \mathcal{X}}$, 
and proposal kernel $q(\cdot\mid\cdot)$. Suppose that $\pi_{C,z}(z) = \pi(z)$ for all $z \in \mathcal{X}$.
Then, the acceptance probability for DA is given by
\begin{align*}
  \alpha(y\mid x) = \min\left\{1, \frac{\min\left\{\pi(y)q(x\mid y), \pi_{C,y}(x)q(y\mid x)\right\}}{\min\left\{\pi(x)q(y\mid x), \pi_{C,x}(y)q(x\mid y)\right\}}\right\}
\end{align*}
for $x \neq y$.
\end{lemma}
\begin{proof}
Fix $x \neq y$. 
Recall that the effective proposal in DA is given by 
\[
  q_C(y\mid x) = q(y\mid x)\min\left\{1, \frac{\pi_{C,x}(y)q(x\mid y)}{\pi_{C,x}(x)q(y\mid x)}\right\}
\]
Then, the outer acceptance probability is given by
\begin{align*}
  \alpha(y\mid x) 
  &= \min\left\{1, \frac{\pi(y)q_C(x\mid y)}{\pi(x)q_C(y\mid x)}\right\} \\
  &= \min\left\{1, \frac{\pi(y)q(x\mid y)\min\left\{1, \frac{\pi_{C,y}(x)q(y\mid x)}{\pi_{C,y}(y)q(x\mid y)}\right\}}{\pi(x)q(y\mid x)\min\left\{1, \frac{\pi_{C,x}(y)q(x\mid y)}{\pi_{C,x}(x)q(y\mid x)}\right\}}\right\} \\
  &= \min\left\{1, \frac{\min\left\{\pi(y)q(x\mid y), \frac{\pi(y)\pi_{C,y}(x)q(y\mid x)}{\pi_{C,y}(y)}\right\}}{\min\left\{\pi(x)q(y\mid x), \frac{\pi(x)\pi_{C,x}(y)q(x\mid y)}{\pi_{C,x}(x)}\right\}}\right\}\\
  &= \min\left\{1, \frac{\min\left\{\pi(y)q(x\mid y), \pi_{C,y}(x)q(y\mid x)\right\}}{\min\left\{\pi(x)q(y\mid x), \pi_{C,x}(y)q(x\mid y)\right\}}\right\},
\end{align*}
where in the final step we used that $\pi_{C, x}(x) = \pi(x)$ and $\pi_{C, y}(y) = \pi(y)$.

Note that this acceptance probability also works even when the densities are not normalised. 
To see this, recall that $F_x(x) = F(x)$, and so $\mathcal L_x(x) = \mathcal L(x)$. 
This implies that
\begin{align*}
  \pi_{C,x}(z) = \frac{\mathcal L_x(z)p(z)}{Z_x} \quad \text{and} \quad \pi(z) = \frac{\mathcal L(z)p(z)}{Z}
\end{align*}
for all $z \in \mathcal X$, where $Z_x = \int_\mathcal X \mathcal L_x(z)p(z) dz$ and $Z = \int_\mathcal X \mathcal L(z)p(z) dz$.
This gives
\begin{align*}
  \frac{\pi(x)\pi_{C,x}(y)}{\pi_{C,x}(x)} = \frac{\mathcal L(y)p(y)}{Z} = \frac{1}{Z} \tilde{\pi}_{C,x}(y),
\end{align*}
where $\tilde{\pi}_{C,x}(y) := \mathcal L(y)p(y)$ is the unnormalized coarse density. 
By the same argument, one can show that this also holds for $\tilde{\pi}_{C,x}$.
As such, we have that 
\begin{align*}
  \frac{\pi(x)\pi_{C,x}(y)}{\pi_{C,x}(x)} = K\cdot \tilde{\pi}_{C,x}(y) 
  \quad \text{and} \quad
  \frac{\pi(y)\pi_{C,y}(x)}{\pi_{C,y}(y)} = K\cdot \tilde{\pi}_{C,y}(x),
\end{align*}
where $K = \frac{1}{Z}$ is a constant. Then, with $\tilde(\cdot) := K \cdot \pi(\cdot)$, 
the acceptance probability becomes
\begin{align*}
  \alpha(y\mid x) 
  &= \min\left\{1, \frac{\min\left\{\pi(y)q(x\mid y), \frac{\pi(y)\pi_{C,y}(x)q(y\mid x)}{\pi_{C,y}(y)}\right\}}{\min\left\{\pi(x)q(y\mid x), \frac{\pi(x)\pi_{C,x}(y)q(x\mid y)}{\pi_{C,x}(x)}\right\}}\right\}\\
  &= \min\left\{1, \frac{\min\left\{K \cdot \tilde{\pi}(y)q(x\mid y), K \cdot \tilde{\pi}_{C,y}(x)q(y\mid x)\right\}}{\min\left\{K \cdot \tilde{\pi}(x)q(y\mid x), K \cdot \tilde{\pi}_{C,x}(y)q(x\mid y)\right\}}\right\}\\
  &= \min\left\{1, \frac{\min\left\{\tilde{\pi}(y)q(x\mid y), \tilde{\pi}_{C,y}(x)q(y\mid x)\right\}}{\min\left\{\tilde{\pi}(x)q(y\mid x), \tilde{\pi}_{C,x}(y)q(x\mid y)\right\}}\right\}.
\end{align*}
\end{proof}

\begin{lemma}[Acceptance probability for DA with state-independent coarse density]
Fix a target density $\pi(\cdot)$, state-independent coarse density $\pi_{C}(\cdot)$, 
and proposal kernel $q(\cdot\mid\cdot)$.
Then, the acceptance probability for DA is given by
\begin{align*}
  \alpha(y\mid x) = \min\left\{1, \frac{\pi(y)\pi_C(x)}{\pi(x)\pi_C(y)}\right\}
\end{align*}
for $x \neq y$.
\end{lemma}
\begin{proof}
Fix $x \neq y$. 
Recall that the effective proposal in DA is given by 
\[
  q_C(y\mid x) = q(y\mid x)\min\left\{1, \frac{\pi_{C}(y)q(x\mid y)}{\pi_{C}(x)q(y\mid x)}\right\}.
\]
Then, the outer acceptance probability is given by
\begin{align*}
  \alpha(y\mid x) 
  &= \min\left\{1, \frac{\pi(y)q_C(x\mid y)}{\pi(x)q_C(y\mid x)}\right\} \\
  &= \min\left\{1, \frac{\pi(y)q(x\mid y)\min\left\{1, \frac{\pi_{C}(x)q(y\mid x)}{\pi_{C}(y)q(x\mid y)}\right\}}{\pi(x)q(y\mid x)\min\left\{1, \frac{\pi_{C}(y)q(x\mid y)}{\pi_{C}(x)q(y\mid x)}\right\}}\right\} \\
  &= \min\left\{1, \frac{\pi(y)\frac{1}{\pi_{C}(y)}\min\left\{\pi_{C}(y)q(x\mid y), \pi_{C}(x)q(y\mid x)\right\}}{\pi(x)\frac{1}{\pi_{C}(x)}\min\left\{\pi_{C}(x)q(y\mid x), \pi_{C}(y)q(x\mid y)\right\}}\right\} \\
  &= \min\left\{1, \frac{\pi(y)\pi_C(x)}{\pi(x)\pi_C(y)}\right\}.
\end{align*}
\end{proof}

Note that DA is a special case of IRM with $n=1$. 
As such, let us fix $n>1$, and attempt to Peskun order DA and IRM. 

\begin{proposition}[Comparison of DA and state-independent IRM]
Fix a target density $\pi(\cdot)$, a state-independent coarse density $\pi_C(\cdot)$,
and a proposal kernel $q(\cdot\mid\cdot)$. Let $Q$ be the Metropolis--Hastings kernel
targeting $\pi_C(\cdot)$ with proposal $q(\cdot\mid\cdot)$.

Let $P_{\mathrm{DA}}$ denote the DA kernel and let
$P_{\mathrm{IRM},n}$ denote the IRM kernel with inner-chain length
$n\geq 1$. Then, for $x\neq y$,
\[
P_{\mathrm{DA}}(x,dy)
=
Q(x,dy)\beta(x,y),
\]
and
\[
P_{\mathrm{IRM},n}(x,dy)
=
Q^n(x,dy)\beta(x,y),
\]
where
\[
\beta(x,y)
=
\min\left\{
1,
\frac{\pi(y)\pi_C(x)}
      {\pi(x)\pi_C(y)}
\right\}.
\]

Consequently,
\[
P_{\mathrm{IRM},n}\succeq_{\mathrm P}P_{\mathrm{DA}}
\]
if and only if, for every $x$ and every measurable set $A$ with
$x\notin A$,
\[
\int_A \beta(x,y)Q^n(x,dy)
\geq
\int_A \beta(x,y)Q(x,dy).
\]
Similarly,
\[
P_{\mathrm{DA}}\succeq_{\mathrm P}P_{\mathrm{IRM},n}
\]
if and only if the reverse inequality holds.

In general, neither Peskun ordering holds.
\end{proposition}

\begin{proof}
Since $Q$ is a Metropolis--Hastings kernel targeting $\pi_C$, it
satisfies detailed balance with respect to $\pi_C$:
\[
\pi_C(x)Q(x,dy)
=
\pi_C(y)Q(y,dx).
\]
Every power $Q^n$ is also reversible with respect to $\pi_C$, and hence
\[
\pi_C(x)Q^n(x,dy)
=
\pi_C(y)Q^n(y,dx).
\]
Therefore, for $x\neq y$,
\[
\frac{Q^n(x\mid y)}{Q^n(y\mid x)}
=
\frac{\pi_C(x)}{\pi_C(y)}.
\]

The IRM outer acceptance probability is consequently
\[
\begin{aligned}
\alpha_n(x,y)
&=
\min\left\{
1,
\frac{\pi(y)Q^n(x\mid y)}
      {\pi(x)Q^n(y\mid x)}
\right\} \\
&=
\min\left\{
1,
\frac{\pi(y)\pi_C(x)}
      {\pi(x)\pi_C(y)}
\right\}
=
\beta(x,y).
\end{aligned}
\]
This expression is independent of $n$. In particular, DA, which corresponds to $n=1$, has the same outer acceptance
probability.

It follows that
\[
P_{\mathrm{DA}}(x,dy)
=
Q(x,dy)\beta(x,y),
\]
whereas
\[
P_{\mathrm{IRM},n}(x,dy)
=
Q^n(x,dy)\beta(x,y)
\]
away from the diagonal. The stated characterisations now follow
directly from the definition of Peskun ordering.

There is no general ordering between the off-diagonal probabilities
of $Q$ and $Q^n$, so there is no general Peskun ordering between
DA and IRM.
\end{proof}

\begin{example}[IRM can Peskun-dominate DA]
Let
\[
\mathcal X=\{0,1\},
\qquad
\pi=\pi_C=\left(\frac12,\frac12\right),
\]
and take
\[
q(1\mid0)=q(0\mid1)=\frac14,
\qquad
q(0\mid0)=q(1\mid1)=\frac34.
\]
Since $\pi_C$ is uniform and $q$ is symmetric, every inner proposal is
accepted, so $Q=q$. Moreover, since $\pi=\pi_C$, the outer acceptance
probability satisfies
\[
\beta(x,y)=1.
\]

For IRM with $n=2$,
\[
Q^2(1\mid0)
=
Q(0\mid0)Q(1\mid0)
+
Q(1\mid0)Q(1\mid1)
=
2\left(\frac34\right)\left(\frac14\right)
=
\frac38.
\]
Thus,
\[
P_{\mathrm{IRM},2}(0,1)
=
\frac38
>
\frac14
=
P_{\mathrm{DA}}(0,1).
\]
The same inequality holds in the reverse direction, and hence
\[
P_{\mathrm{IRM},2}
\succeq_{\mathrm P}
P_{\mathrm{DA}}.
\]
\end{example}

\begin{example}[DA can Peskun-dominate IRM]
Again let
\[
\mathcal X=\{0,1\},
\qquad
\pi=\pi_C=\left(\frac12,\frac12\right),
\]
but now take
\[
q(1\mid0)=q(0\mid1)=\frac34,
\qquad
q(0\mid0)=q(1\mid1)=\frac14.
\]
As before, $Q=q$ and $\beta(x,y)=1$. However,
\[
Q^2(1\mid0)
=
2\left(\frac14\right)\left(\frac34\right)
=
\frac38.
\]
Therefore,
\[
P_{\mathrm{DA}}(0,1)
=
\frac34
>
\frac38
=
P_{\mathrm{IRM},2}(0,1),
\]
and similarly in the reverse direction. Hence
\[
P_{\mathrm{DA}}
\succeq_{\mathrm P}
P_{\mathrm{IRM},2}.
\]
\end{example}

The preceding examples show that, even with the same proposal kernel
$q$ and the same state-independent coarse density $\pi_C$, there is no
universal Peskun ordering between DA and IRM.
The common outer acceptance probability means that the comparison is
entirely determined by the off-diagonal behaviour of $Q$ and $Q^n$.

With that said, if the state-independent coarse density $\pi_C(\cdot)$ 
is well-chosen (i.e., approximates the target density $\pi(\cdot)$ well enough, in some sense), 
then the idea is that after $n$ steps, the sub chain will end up in a different mode.
That is, in a region far from the current state that has high probability under 
the target density $\pi(\cdot)$. 
In this case, $Q^n$ will have much more probability mass on the off-diagonal than $Q$ 
(especially in regions where $\beta(x, \cdot)$ is not low?), and so 
IRM should Peskun-dominate DA.

Since there is no general Peskun ordering between DA and IRM with a state-independent coarse density and the same proposal, 
the same is true for DA and IRM with state-dependent coarse densities and the same proposal. 


\section{Why is it difficult?}
\hl{TODO.}
Explain why it is difficult to compute the exact MH acceptance probability. 


\section{Initial attempts}

\begin{algorithm}[H]
  \caption{Rat-Mice 1.0 (RM1)}
  \label{alg:rm1}
  \textbf{function:} $[x^1,\ldots,x^N] =
  \mathrm{RM1}\bigl(\pi(\cdot), \{\pi_{C,x}(\cdot)\}_{x \in \mathcal{X}}, q(\cdot\mid\cdot), n, x^0, N\bigr)$
  
  \textbf{input:} target density $\pi(\cdot)$, state-dependent surrogate (coarse) densities $\{\pi_{C,x}(\cdot)\}_{x \in \mathcal{X}}$, 
  proposal kernel $q(\cdot\mid\cdot)$, subchain length $n$, initial state $x^0$, number of steps $N$
  
  \textbf{output:} ordered list of states $[x^1,\ldots,x^N]$
  
  \begin{algorithmic}[1]
  \For{$j = 0$ to $N-1$}
      \State Starting at $x^j$, generate subchain of length $n$ using MH Alg. 1 to target $\pi_{C, x_j}(\cdot)$:
      \[
          z = \mathrm{MH}\bigl(\pi_{C, x_j}(\cdot), q(\cdot\mid\cdot), x^j, n\bigr).
      \]
      \State Accept the state $z$ as our proposal, i.e. set $y = z$, with probability
      \[
          \alpha_{x_j}(z \mid x^j)
          =
          \min\left\{
              1,
              \frac{\pi(z)\pi_{C, x_j}(x^j)}
                   {\pi(x^j)\pi_{C, x_j}(z)}
          \right\}.
      \]
      \State Otherwise reject $z$ and set $y = x^j$.
      \State Accept the proposal $y$ as the next sample, i.e. set $x^{j+1} = y$, with probability
      \[
          \alpha(y \mid x^j)
          = 1.
      \]
  \EndFor
  \end{algorithmic}
\end{algorithm}


\begin{algorithm}[H]
  \caption{Pathwise Rat-Mice (PRM)}
  \label{alg:prm}

  \textbf{function:}
  $[x^1,\ldots,x^N]
  =
  \mathrm{PRM}
  \bigl(\pi(\cdot), \{\pi_{C,x}(\cdot)\}_{x\in\mathcal X},
  q(\cdot\mid\cdot), n, x^0, N\bigr)$

  \textbf{input:} target density $\pi(\cdot)$, state-dependent surrogate (coarse)
  densities $\{\pi_{C,x}(\cdot)\}_{x\in\mathcal X}$, proposal density
  $q(\cdot\mid\cdot)$, subchain length $n$, initial state $x^0$, number of steps $N$

  \textbf{output:} ordered list of states $[x^1,\ldots,x^N]$

  \begin{algorithmic}[1]
  \For{$j=0$ to $N-1$}
      \State Set $z_0=x^j$ and set $\ell_F=0$.
      \For{$k=0$ to $n-1$}
          \State Propose $u_k\sim q(\cdot\mid z_k)$.
          \State Compute the MH acceptance probability targeting $\pi_{C,x^j}$:
          \[
              \beta_{x^j}(u_k\mid z_k)
              =
              \min\left\{
              1,
              \frac{
                  \pi_{C,x^j}(u_k)q(z_k\mid u_k)
              }{
                  \pi_{C,x^j}(z_k)q(u_k\mid z_k)
              }
              \right\}.
          \]
          \State Draw $r_k\sim \mathrm{Bernoulli}(\beta_{x^j}(u_k\mid z_k))$.
          \If{$r_k=1$}
              \State Set $z_{k+1}=u_k$.
              \State Update $\ell_F \leftarrow \ell_F + \log q(u_k\mid z_k) + \log \beta_{x^j}(u_k\mid z_k)$.
          \Else
              \State Set $z_{k+1}=z_k$.
              \State Update $\ell_F \leftarrow \ell_F + \log q(u_k\mid z_k) + \log\left(1-\beta_{x^j}(u_k\mid z_k)\right)$.
          \EndIf
      \EndFor
      \State Set $y=z_n$ and $\ell_R=0$.
      \For{$k=n-1$ down to $0$}
          \If{$r_k=1$}
              \State Compute the reverse MH acceptance probability targeting $\pi_{C,y}$:
              \[
                  \beta_y(z_k\mid z_{k+1})
                  =
                  \min\left\{
                  1,
                  \frac{
                      \pi_{C,y}(z_k)q(z_{k+1}\mid z_k)
                  }{
                      \pi_{C,y}(z_{k+1})q(z_k\mid z_{k+1})
                  }
                  \right\}.
              \]
              \State Update $\ell_R \leftarrow \ell_R + \log q(z_k\mid z_{k+1}) + \log \beta_y(z_k\mid z_{k+1})$.
          \Else
              \State Compute the reverse MH rejection probability targeting $\pi_{C,y}$:
              \[
                  \beta_y(u_k\mid z_k)
                  =
                  \min\left\{
                  1,
                  \frac{
                      \pi_{C,y}(u_k)q(z_k\mid u_k)
                  }{
                      \pi_{C,y}(z_k)q(u_k\mid z_k)
                  }
                  \right\}.
              \]
              \State Update $\ell_R \leftarrow \ell_R + \log q(u_k\mid z_k) + \log\left(1-\beta_y(u_k\mid z_k)\right)$.
          \EndIf
      \EndFor
      \State Accept $y$ as the next sample, i.e. set $x^{j+1}=y$, with probability
      \[
          \alpha(y\mid x^j)
          =
          \min\left\{
          1,
          \exp\left[
              \log\pi(y)+\ell_R-\log\pi(x^j)-\ell_F
          \right]
          \right\}.
      \]
      \State Otherwise reject $y$ and set $x^{j+1}=x^j$.
  \EndFor
  \end{algorithmic}
\end{algorithm}


\section{RM1 is not in detailed balance}

\begin{lemma}
\label{lem:rm1-detailed-balance}
Fix a target density $\pi(\cdot)$, state-dependent densities $\{\pi_{C,x}(\cdot)\}_{x \in \mathcal{X}}$, 
proposal kernel $q(\cdot\mid\cdot)$, and subchain length $n$. 
Then, the RM1 algorithm does not in general satisfy the detailed balance condition with respect to $\pi$.
\end{lemma}
\begin{proof}
Let $P$ denote the transition kernel of the Markov chain generated by the RM1 algorithm.
Let $p(x, y)$ denote the density of the transition kernel $P$, i.e.,
\[
P(x, A) = \int_A p(x, y) dy, \quad \forall A\in\mathcal B(\mathcal X).
\]
We aim to show that 
\[
  p(x, y) \pi(x) = p(y, x) \pi(y), \quad \forall x,y\in\mathcal X.
\]
To this end, fix $x \in \mathcal X$. 
Let $q_x(b|a)$ denote the transition kernel obtained by 
running the MH algorithm for one step, starting from $a$ and targeting $\pi_{C,x}(\cdot)$ with proposal kernel $q(\cdot\mid\cdot)$.
That is,
\[
q_x(b|a) = q(b|a) \min\left\{1, \frac{\pi_{C,x}(b)q(a|b)}{\pi_{C,x}(a)q(b|a)}\right\} + \delta_a(b)\left(1 - \int_\mathcal X q(c|a) \min\left\{1, \frac{\pi_{C,x}(c)q(a|c)}{\pi_{C,x}(a)q(c|a)}\right\} dc\right).
\]
By design of the MH algorithm, $q_x(\cdot|\cdot)$ satisfies the detailed balance condition with respect to $\pi_{C,x}(\cdot)$. 
It then follows that $q_x^n(\cdot|\cdot)$ also satisfies detailed balance with respect to $\pi_{C,x}(\cdot)$ 
(see Lemma 2 in Lykkegaard et al. 2023).
That is, 
\begin{align}
  q_x^n(a|b) \pi_{C,x}(b) &= q_x^n(b|a) \pi_{C,x}(a), \quad \forall a,b\in\mathcal X.\nonumber
\end{align}
For all $a,b\in\mathcal X$, where $\pi_{C,x}(b), q_x^n(b|a) > 0$, this gives,
\begin{align}
  \frac{q_x^n(a|b)}{q_x^n(b|a)} &= \frac{\pi_{C,x}(a)}{\pi_{C,x}(b)}. 
\end{align}
Using this, steps 2 to 4 in RM1 
define a proposal kernel $q^*(\cdot|\cdot)$ given by
\begin{align}
  q^*(y|x) 
  &= q_x^n(y|x)\alpha_{x}(y|x) + \delta_x(y)\left(1 - \int_\mathcal X q_x^n(z|x) \alpha_{x}(z|x) dz\right) \nonumber\\
  &= q_x^n(y|x)\min\left\{1, \frac{\pi(y)\pi_{C,x}(x)}{\pi(x)\pi_{C,x}(y)}\right\} + \delta_x(y)\left(1 - \int_\mathcal X q_x^n(z|x) \min\left\{1, \frac{\pi(z)\pi_{C,x}(x)}{\pi(x)\pi_{C,x}(z)}\right\} dz\right) \nonumber\\
  &= q_x^n(y|x)\min\left\{1, \frac{\pi(y)q_x^n(x|y)}{\pi(x)q_x^n(y|x)}\right\} + \delta_x(y)\left(1 - \int_\mathcal X q_x^n(z|x) \min\left\{1, \frac{\pi(z)q_x^n(x|z)}{\pi(x)q_x^n(z|x)}\right\} dz\right).
  \nonumber
\end{align}
Now, we have that for $y \neq x$, 
\begin{align}
  \pi(x) q^*(y|x) 
  = \pi(x) q_x^n(y|x)\min\left\{1, \frac{\pi(y)q_x^n(x|y)}{\pi(x)q_x^n(y|x)}\right\}
  = \min\left\{\pi(x)q_x^n(y|x),\pi(y)q_x^n(x|y)\right\}. \label{eq:detailed-balance-x-y}
\end{align}
On the other hand, 
\begin{align}
  \pi(y) q^*(x|y) 
  = \pi(y) q_y^n(x|y)\min\left\{1, \frac{\pi(x)q_y^n(y|x)}{\pi(y)q_y^n(x|y)}\right\}
  = \min\left\{\pi(y)q_y^n(x|y),\pi(x)q_y^n(y|x)\right\}. \label{eq:detailed-balance-y-x}
\end{align}
In general, \eqref{eq:detailed-balance-x-y} and \eqref{eq:detailed-balance-y-x} are not equal. 
Thus, $q^*(\cdot|\cdot)$ does not automatically satisfy detailed balance with respect to $\pi$.
\end{proof}

\begin{remark}
  One can show that if we want to choose $q^*(y|x)$ of the form
  \begin{align}
    q^*(y|x) 
    = q_x^n(y|x)\min\left\{1, K(y|x)\right\}, \label{eq:q-star-form}
  \end{align}
  and enforce detailed balance with respect to $\pi(\cdot)$, then we must have that 
  \begin{align*}
    K(y|x) = \frac{\pi(y)q_y^n(x|y)}{\pi(x)q_x^n(y|x)}.
  \end{align*}
  This is precisely the MH term in the acceptance probability. 
  Thus, if we want to have an MH acceptance probability of the form $\min\left\{1, K(y|x)\right\}$, 
  then our only choice is to compute $q_y^n(x|y)/q_x^n(y|x)$.
\end{remark}


% -------- PRM --------
\section{PRM may produce small acceptance probabilities}

\begin{lemma} 
\label{lem:prm-low-acceptance}
Fix a target density $\pi(\cdot)$, state-dependent densities $\{\pi_{C,x}(\cdot)\}_{x \in \mathcal{X}}$, 
proposal kernel $q(\cdot\mid\cdot)$, and subchain length $n$. 
Then, the PRM algorithm may produce small acceptance probabilities.
\end{lemma}

\begin{proof}
Observe that 
\begin{align*}
  \exp\left(\ell_F\right) 
  &= \prod_{k=0}^{n-1} \begin{cases}
    q(u_k|z_k) \beta_{x^j}(u_k|z_k) & \text{if } r_k = 1 \\
    q(u_k|z_k) (1-\beta_{x^j}(u_k|z_k)) & \text{if } r_k = 0
  \end{cases}.
\end{align*}
On the other hand, 
\begin{align*}
  \exp\left(\ell_R\right) 
  &= \prod_{k=n-1}^{0} \begin{cases}
    q(z_k|z_{k+1}) \beta_y(z_k|z_{k+1}) & \text{if } r_k = 1 \\
    q(u_k|z_{k}) (1-\beta_y(u_k|z_{k})) & \text{if } r_k = 0
  \end{cases}\\
  &= \prod_{k=n-1}^{0} \begin{cases}
    q(z_k|u_k) \beta_y(z_k|u_k) & \text{if } r_k = 1 \\
    q(u_k|z_{k}) (1-\beta_y(u_k|z_{k})) & \text{if } r_k = 0
  \end{cases}.
\end{align*}
Now, if $r_k = 1$, then
\begin{align*}
  \frac{q(z_k|u_k) \beta_y(z_k|u_k)}{q(u_k|z_k) \beta_{x^j}(u_k|z_k)}
  &= \frac{q(z_k|u_k) \min\left\{1, \frac{\pi_{C,y}(z_k)q(u_k|z_k)}{\pi_{C,y}(u_k)q(z_k|u_k)}\right\}}{q(u_k|z_k) \min\left\{1, \frac{\pi_{C,x^j}(u_k)q(z_k|u_k)}{\pi_{C,x^j}(z_k)q(u_k|z_k)}\right\}}.
\end{align*}
In the case where $q(\cdot|\cdot)$ is symmetric, this becomes 
\[
  \frac{\min\left\{1, \frac{\pi_{C,y}(z_k)}{\pi_{C,y}(u_k)}\right\}}{\min\left\{1, \frac{\pi_{C,x^j}(u_k)}{\pi_{C,x^j}(z_k)}\right\}}.
\]
The denominator is probably large (close to or larger than $1$) as with $r_k=1$ we would have accepted the proposal. 
On the other hand, the numerator is probably small since $u_k$ is `closer' to $y$ than $z_k$, 
and $\pi_{C,y}(\cdot)$ is probably larger than $\pi_{C,x^j}(\cdot)$. 
As such, the numerator is probably less than $1$. 
In the case where the denominator is more than $1$,
the full term is less than $1$.
In the case where the denominator is also less than $1$, we have
\[
  \frac{\frac{\pi_{C,y}(z_k)}{\pi_{C,y}(u_k)}}{\frac{\pi_{C,x^j}(u_k)}{\pi_{C,x^j}(z_k)}}
  = \frac{\pi_{C,y}(z_k)}{\pi_{C,y}(u_k)} \cdot \frac{\pi_{C,x^j}(z_k)}{\pi_{C,x^j}(u_k)}.
\]
I suspect this term is probably less than $1$ as with $r_k=1$, $\frac{\pi_{C,x^j}(u_k)}{\pi_{C,x^j}(z_k)}$ 
will be quite close to $1$, and so $\frac{\pi_{C,y}(z_k)}{\pi_{C,y}(u_k)}$ also be not much larger than $1$.
On the other hand, the $\frac{\pi_{C,y}(z_k)}{\pi_{C,y}(u_k)}$ may well be quite a bit smaller than $1$. 

Long story short, we probably have that 
\[
  \frac{q(z_k|u_k) \beta_y(z_k|u_k)}{q(u_k|z_k) \beta_{x^j}(u_k|z_k)} < 1.
\]

Now, if $r_k = 0$, then
\begin{align*}
  \frac{q(u_k|z_k) (1-\beta_y(u_k|z_k))}{q(u_k|z_k) (1-\beta_{x^j}(u_k|z_k))}
  &= \frac{1 - \min\left\{1, \frac{\pi_{C,y}(u_k)q(z_k|u_k)}{\pi_{C,y}(z_k)q(u_k|z_k)}\right\}}{1 - \min\left\{1, \frac{\pi_{C,x^j}(u_k)q(z_k|u_k)}{\pi_{C,x^j}(z_k)q(u_k|z_k)}\right\}}.
\end{align*}
In the case where $q(\cdot|\cdot)$ is symmetric, this becomes 
\[
  \frac{1 - \min\left\{1, \frac{\pi_{C,y}(u_k)}{\pi_{C,y}(z_k)}\right\}}{1 - \min\left\{1, \frac{\pi_{C,x^j}(u_k)}{\pi_{C,x^j}(z_k)}\right\}}.
\]
By definition, since $r_k = 0$, we have that that $\frac{\pi_{C,x^j}(u_k)}{\pi_{C,x^j}(z_k)}$ is less than $1$. 
So, this becomes 
\[
  \frac{1 - \min\left\{1, \frac{\pi_{C,y}(u_k)}{\pi_{C,y}(z_k)}\right\}}{1 - \frac{\pi_{C,x^j}(u_k)}{\pi_{C,x^j}(z_k)}}.
\]
Since we rejected the proposal, $z_k$ is probably `closer' to $y$ than $u_k$, 
and so $\pi_{C,y}(z_k)$ is probably larger than $\pi_{C,y}(u_k)$. This gives
\[
  \frac{1 - \frac{\pi_{C,y}(u_k)}{\pi_{C,y}(z_k)}}{1 - \frac{\pi_{C,x^j}(u_k)}{\pi_{C,x^j}(z_k)}}.
\]
Finally, since we rejected the proposal, the denominator may well be larger than the numerator.
As such, we probably have that
\[
  \frac{q(u_k|z_k) (1-\beta_y(u_k|z_k))}{q(u_k|z_k) (1-\beta_{x^j}(u_k|z_k))} < 1.
\]

Overall, we have that 
\[  
  \frac{\exp\left(\ell_F\right)}{\exp\left(\ell_R\right)}
\]
is a product of $n$ terms, each of which is probably less than $1$. 
So, with more terms, the acceptance probability may geometrically decay. 

Furthermore, if we ever have that 
$r_k=0$ and $\beta_{y}(u_k|z_k) = 1$, 
then the entire acceptance probability is equal to $0$.
\end{proof}


% -------- PRM --------
\section{Mean path-dependent acceptance probability is bounded by MH acceptance probability}

\begin{lemma}
Let $\Omega_{x,y}$ denote the set of auxiliary paths from $x$ to $y$, with
path density $Q_x(\omega)$, and define the marginal proposal density
\[
q_x^n(y\mid x)
=
\int_{\Omega_{x,y}} Q_x(\omega)\,d\omega.
\]
Suppose that $\alpha_x(\omega)\in[0,1]$ is a path-dependent acceptance
probability such that the resulting accepted flows satisfy
\[
\pi(x)
\int_{\Omega_{x,y}}
Q_x(\omega)\alpha_x(\omega)\,d\omega
=
\pi(y)
\int_{\Omega_{y,x}}
Q_y(\omega')\alpha_y(\omega')\,d\omega'.
\]
Then, whenever $q_x^n(y\mid x)>0$, the conditional mean acceptance
probability
\[
\bar\alpha(x,y)
:=
\frac{
\int_{\Omega_{x,y}}
Q_x(\omega)\alpha_x(\omega)\,d\omega
}{
q_x^n(y\mid x)
}
\]
satisfies
\[
\bar\alpha(x,y)
\le
\min\left\{
1,\,
\frac{\pi(y)q_y^n(x\mid y)}
      {\pi(x)q_x^n(y\mid x)}
\right\}.
\]
Thus, for a fixed marginal proposal density, no reversible
path-dependent acceptance rule has a larger mean acceptance probability
than the marginal Metropolis--Hastings rule.
\end{lemma}

\begin{proof}
Define the accepted probability flow from $x$ to $y$ by
\[
A(x,y)
:=
\int_{\Omega_{x,y}}
Q_x(\omega)\alpha_x(\omega)\,d\omega.
\]

Since $\alpha_x(\omega)\le 1$ for every path $\omega$,
\[
A(x,y)
\le
\int_{\Omega_{x,y}}Q_x(\omega)\,d\omega
=
q_x^n(y\mid x).
\]

By detailed balance of the accepted endpoint transitions,
\[
\pi(x)A(x,y)=\pi(y)A(y,x),
\]
and hence
\[
A(x,y)=\frac{\pi(y)}{\pi(x)}A(y,x).
\]
Applying the same bound to the reverse move gives
\[
A(y,x)
=
\int_{\Omega_{y,x}}
Q_y(\omega')\alpha_y(\omega')\,d\omega'
\le
\int_{\Omega_{y,x}}Q_y(\omega')\,d\omega'
=
q_y^n(x\mid y).
\]
Therefore,
\[
A(x,y)
\le
\frac{\pi(y)}{\pi(x)}q_y^n(x\mid y).
\]

Combining the two bounds yields
\[
A(x,y)
\le
\min\left\{
q_x^n(y\mid x),
\frac{\pi(y)}{\pi(x)}q_y^n(x\mid y)
\right\}.
\]
Dividing by $q_x^n(y\mid x)>0$, we obtain
\[
\bar\alpha(x,y)
=
\frac{A(x,y)}{q_x^n(y\mid x)}
\le
\min\left\{
1,
\frac{\pi(y)q_y^n(x\mid y)}
      {\pi(x)q_x^n(y\mid x)}
\right\},
\]
which is the Metropolis--Hastings acceptance probability for the
marginal endpoint proposal.
\end{proof}


% \section{Definitions and Useful Results}

% We list the following definitions before stating 
% Theorem \ref{thm:ergodic}, which we will use to prove the ergodicity of the RM algorithm.
% This theorem is a rewording of Theorem 6.51 in Monte Carlo Statistical Methods 
% by Robert and Casella.

% \begin{definition}[Transition kernel]
% A \emph{transition kernel} on $(\mathcal X,\mathcal B(\mathcal X))$ is a map
% \[
% P:\mathcal X\times \mathcal B(\mathcal X) \to [0,1]
% \]
% such that:
% \begin{enumerate}
%     \item for each $x\in\mathcal X$, $P(x,\cdot)$ is a probability measure on $(\mathcal X,\mathcal B(\mathcal X))$;
%     \item for each $A\in\mathcal B(\mathcal X)$, $P(\cdot,A)$ is $\mathcal B(\mathcal X)$-measurable.
% \end{enumerate}
% Here $P(x,A)$ is interpreted as the probability that the chain moves from state $x$ into the set $A$ in one step.
% \end{definition}

% \begin{definition}[$n$-step transition kernel]
% For $n\geq 1$, the $n$-step transition kernel $P^n$ is defined recursively by
% \[
% P^{n+1}(x,A)
% =
% \int_{\mathcal X} P^n(y,A) P(x,dy),
% \qquad \forall x\in\mathcal X, \forall A\in\mathcal B(\mathcal X).
% \]
% Thus $P^n(x,A)$ is the probability that the Markov chain is in $A$ after $n$ steps when started from $x$.
% \end{definition}

% \begin{definition}[Invariant distribution]
% A probability measure $\pi$ on $(\mathcal X,\mathcal B(\mathcal X))$ is called an \emph{invariant distribution} for $P$ if
% \[
% \pi(A)
% =
% \int_{\mathcal X} P(x,A)\,\pi(dx),
% \qquad \forall A\in\mathcal B(\mathcal X).
% \]
% \end{definition}

% \begin{definition}[$\pi$-irreducibility]
% Let $\pi$ be a probability measure on $(\mathcal X,\mathcal B(\mathcal X))$. A transition kernel $P$ is called
% \emph{$\pi$-irreducible} if, for every measurable set $A\in\mathcal B(\mathcal X)$ with $\pi(A)>0$, there exists some $n\geq 1$ such that,
% \[
% P^n(x,A)>0, \quad \forall x\in\mathcal X.
% \]
% \end{definition}

% \begin{definition}[Harris recurrence]
% Suppose that $P$ is $\pi$-irreducible. A Markov chain 
% with transition kernel $P$ 
% is called \emph{Harris recurrent} if, for every measurable set $A\in\mathcal B(\mathcal X)$ with $\pi(A)>0$,
% \[
% \mathbb P_x(\tau_A<\infty)=1,
% \qquad \forall x\in\mathcal X,
% \]
% where
% \[
% \tau_A=\inf\{n\geq 1:X_n\in A\}.
% \]
% \end{definition}

% \hl{Note that this definition of aperiodicity is different from the one in Monte Carlo Statistical Methods.}
% \begin{definition}[Aperiodicity]
% Let $P$ be a $\pi$-irreducible transition kernel on$(\mathcal X,\mathcal B(\mathcal X))$. The chain is called
% \emph{aperiodic} if there is no integer $d\geq 2$ and no collection of disjoint measurable sets
% \[
% D_0,D_1,\ldots,D_{d-1}\in\mathcal B(\mathcal X),
% \qquad \pi(D_i)>0,
% \]
% such that
% \[
% P(x,D_{i+1})=1
% \]
% for $\pi$-almost every $x\in D_i$, where the index is taken modulo $d$.
% \end{definition}

% \begin{definition}[Ergodicity]
% A transition kernel $P$ with invariant distribution $\pi$ is called \emph{ergodic} if
% \[
% \norm{P^n(x,\cdot)-\pi(\cdot)}_{TV}\to 0
% \]
% as $n\to\infty$, for all $x\in\mathcal X$.
% \end{definition}

% \begin{theorem}
% \label{thm:ergodic}
% Suppose that \(P\) is a transition kernel on \((\mathcal X,\mathcal B(\mathcal X))\)
% with invariant distribution \(\pi\). If \(P\) is \(\pi\)-irreducible, aperiodic, and Harris recurrent, then \(P\) is ergodic. That is,
% \[
%     \norm{P^n(x,\cdot)-\pi(\cdot)}_{\text{TV}}  = \sup_{A\in\mathcal B(\mathcal X)} |P^n(x,A) - \pi(A)| \to 0
% \]
% as \(n \to \infty\), for all \(x \in \mathcal X\).
% \end{theorem}

% \begin{definition}[Detailed balance]
% Let $P$ be a transition kernel and $\pi$ be a probability measure on $(\mathcal X,\mathcal B(\mathcal X))$.
% We say that $P$ satisfies the \emph{detailed balance condition} with respect to
% $\pi$ if
% \[
%     \int_A P(x,B)\,\pi(dx)
%     =
%     \int_B P(x,A)\,\pi(dx),
%     \qquad
%     \forall A,B\in\mathcal B(\mathcal X).
% \]
% \end{definition}

% \begin{lemma}
% \label{lem:detailed-balance-invariant}
% If a transition kernel $P$ satisfies the detailed balance condition with respect
% to a probability measure $\pi$, then $\pi$ is an invariant distribution for $P$.
% \end{lemma}


% -------- REFERENCES ---------
% \bibliographystyle{ieeetr}
% \bibliography{references}


\end{document}